\FloatBarrier
\section{یکسانی نتایج \lr{rollout} و \lr{attention flow} بین شرکت‌کننده‌ها}
\label{sec:consistency}

\paragraph{چرا یکسانی مهم است.} هر ۱۲ نفر یک کار را با یک دست انجام داده‌اند، و برای هر کدام یک مدل با معماری یکسان جداگانه آموزش دیده است. اگر نقشه‌های تفسیر چیزی واقعی درباره‌ی رابطه‌ی مغز و عضله نشان دهند، باید در افراد مختلف کم‌وبیش شبیه باشند. اگر نقشه در هر نفر شکل کاملاً متفاوتی داشته باشد، یا از نویز آمده یا از جزئیات تصادفی آموزش همان مدل. در هر دو حالت نمی‌توان از آن نتیجه‌ی علمی گرفت.

این بخش یکسانی را برای همه‌ی نقشه‌های ذخیره‌شده اندازه می‌گیرد. هیچ مدلی دوباره اجرا نشده است؛ همه‌ی اعداد از جدول‌های بخش~\ref{sec:interp} و~\ref{sec:rollout-flow} می‌آیند (اسکریپت \lr{v2-report/analysis/consistency.py}).

\subsection{معیارها}

برای هر نوع نقشه یک ماتریس $X\in\mathbb{R}^{12\times n}$ ساخته می‌شود. هر سطر نقشه‌ی یک شرکت‌کننده است (میانگین روی رخدادها) و $n$ تعداد خانه‌ها: ۳۲ الکترود، یا نقاط زمانی، یا ۲۰ درایه‌ی غیرقطری \lr{muscle graph}. پنج معیار، هر کدام از یک زاویه:

\begin{itemize}
\item \textbf{\lr{pairwise $\rho$}: دو نفر چقدر شبیه‌اند؟} نقشه‌ی هر دو نفر با همبستگی اسپیرمن مقایسه و میانگین روی هر ۶۶ جفت گرفته می‌شود (برای منحنی‌های زمانی همبستگی پیرسون). اسپیرمن فقط ترتیب را می‌سنجد: آیا الکترودی که در یک نفر مهم است، در نفر دیگر هم مهم است؟
\item \textbf{\lr{leave-one-out (LOO) $\rho$}: هر نفر چقدر شبیه بقیه‌ی گروه است؟} نقشه‌ی هر نفر با میانگین ۱۱ نفر دیگر مقایسه و نتیجه روی ۱۲ نفر میانگین گرفته می‌شود. میانگین گروه نویز کمتری از یک نفر دارد، پس این عدد معمولاً از \lr{pairwise} بزرگ‌تر است.
\item \textbf{\lr{split-half reliability}: اگر آزمایش را با ۱۲ نفر دیگر تکرار کنیم، همین نقشه‌ی گروهی را می‌گیریم؟} شرکت‌کننده‌ها به‌طور تصادفی به دو گروه ۶ نفری تقسیم می‌شوند و همبستگی میانگین دو گروه حساب می‌شود. چون هر نیمه فقط ۶ نفر است، عدد با فرمول اسپیرمن-براون، $2r/(1+r)$، به اندازه‌ی ۱۲ نفر تصحیح می‌شود. این کار ۵۰۰ بار تکرار و میانگین گرفته می‌شود.
\item \textbf{\lr{Kendall's W}: ۱۲ نفر چقدر در رتبه‌بندی توافق دارند؟} هر شرکت‌کننده مثل یک داور خانه‌ها را رتبه‌بندی می‌کند. $W$ بین ۰ (هیچ توافقی) و ۱ (همه رتبه‌بندی یکسان) است. معناداری با آزمون $\chi^2=m(n-1)W$ سنجیده می‌شود:
\begin{equation}
W=\frac{12\sum_{j=1}^{n}\big(R_j-\bar R\big)^2}{m^2(n^3-n)},\qquad R_j=\sum_{i=1}^{m}\mathrm{rank}_i(j).
\end{equation}
\item \textbf{\lr{null}: اگر هیچ الگوی مشترکی نبود، چه عددی می‌گرفتیم؟} برچسب الکترودها در هر شرکت‌کننده به‌طور مستقل به‌هم ریخته و \lr{pairwise $\rho$} دوباره حساب می‌شود (برای منحنی‌های زمانی، منحنی هر نفر به‌طور دایره‌ای در زمان جابه‌جا می‌شود). این کار ۵۰۰ بار تکرار می‌شود. \lr{null 95\%} صدک ۹۵ این توزیع است. $p$ سهم دفعاتی است که عدد به‌هم‌ریخته از عدد واقعی بزرگ‌تر شده؛ کمترین $p$ ممکن \lr{۰٫۰۰۲} است.
\end{itemize}

\subsection{نقشه‌ی الکترودها}

\begin{table}[htbp]
\centering
\footnotesize
\caption{یکسانی نقشه‌ی الکترودها بین ۱۲ شرکت‌کننده. \lr{pairwise $\rho$}: میانگین همبستگی اسپیرمن روی ۶۶ جفت؛ \lr{LOO}: هر نفر در برابر میانگین ۱۱ نفر دیگر؛ \lr{split-half}: همبستگی میانگین دو نیمه‌ی تصادفی با تصحیح اسپیرمن-براون؛ \lr{null 95\%} و \lr{p}: در برابر جابه‌جایی برچسب الکترودها در هر شرکت‌کننده.}
\label{tab:cons-elec}
\begin{tabular}{lcccccc}
\toprule
نقشه & \lr{pairwise $\rho$} & \lr{LOO $\rho$} & \lr{split-half} & \lr{Kendall's W} & \lr{null 95\%} & \lr{p} \\
\midrule
\lr{E, first-touch, rollout} & \lr{۰٫۲۱} & \lr{۰٫۳۶} & \lr{۰٫۷۰} & \lr{۰٫۲۸} & \lr{۰٫۰۴} & \lr{۰٫۰۰۲} \\
\lr{E, first-touch, gradient-weighted rollout} & \lr{۰٫۲۲} & \lr{۰٫۳۶} & \lr{۰٫۶۸} & \lr{۰٫۲۸} & \lr{۰٫۰۴} & \lr{۰٫۰۰۲} \\
\lr{E, first-touch, attention flow} & \lr{۰٫۲۳} & \lr{۰٫۳۷} & \lr{۰٫۷۰} & \lr{۰٫۲۹} & \lr{۰٫۰۴} & \lr{۰٫۰۰۲} \\
\lr{E, lift-off, rollout} & \lr{۰٫۲۲} & \lr{۰٫۳۶} & \lr{۰٫۶۹} & \lr{۰٫۲۸} & \lr{۰٫۰۴} & \lr{۰٫۰۰۲} \\
\lr{E, lift-off, gradient-weighted rollout} & \lr{۰٫۲۱} & \lr{۰٫۳۶} & \lr{۰٫۶۸} & \lr{۰٫۲۸} & \lr{۰٫۰۳} & \lr{۰٫۰۰۲} \\
\lr{E, lift-off, attention flow} & \lr{۰٫۲۰} & \lr{۰٫۳۴} & \lr{۰٫۶۷} & \lr{۰٫۲۶} & \lr{۰٫۰۳} & \lr{۰٫۰۰۲} \\
\lr{F-eeg, first-touch, rollout} & \lr{۰٫۲۲} & \lr{۰٫۳۵} & \lr{۰٫۶۷} & \lr{۰٫۲۸} & \lr{۰٫۰۴} & \lr{۰٫۰۰۲} \\
\lr{F-eeg, first-touch, gradient-weighted rollout} & \lr{۰٫۱۹} & \lr{۰٫۳۱} & \lr{۰٫۶۲} & \lr{۰٫۲۶} & \lr{۰٫۰۴} & \lr{۰٫۰۰۲} \\
\lr{F-eeg, first-touch, attention flow} & \lr{۰٫۲۴} & \lr{۰٫۴۰} & \lr{۰٫۷۵} & \lr{۰٫۳۱} & \lr{۰٫۰۴} & \lr{۰٫۰۰۲} \\
\lr{F-eeg, lift-off, rollout} & \lr{۰٫۲۲} & \lr{۰٫۳۶} & \lr{۰٫۶۸} & \lr{۰٫۲۸} & \lr{۰٫۰۴} & \lr{۰٫۰۰۲} \\
\lr{F-eeg, lift-off, gradient-weighted rollout} & \lr{۰٫۱۹} & \lr{۰٫۳۱} & \lr{۰٫۶۳} & \lr{۰٫۲۶} & \lr{۰٫۰۴} & \lr{۰٫۰۰۲} \\
\lr{F-eeg, lift-off, attention flow} & \lr{۰٫۲۳} & \lr{۰٫۳۹} & \lr{۰٫۷۴} & \lr{۰٫۳۰} & \lr{۰٫۰۴} & \lr{۰٫۰۰۲} \\
\bottomrule
\end{tabular}
\end{table}


جدول~\ref{tab:cons-elec} نتیجه‌ی اصلی است و دو پیام دارد.

\paragraph{یکسانی واقعی است.} هر ۱۲ نقشه‌ی الکترود (دو خروجی، دو رخداد، سه روش) بین شرکت‌کننده‌ها بسیار بیشتر از \lr{null} شبیه‌اند؛ $p$ در همه کمترین مقدار ممکن است.

\paragraph{اما اندازه‌اش متوسط است.} \lr{pairwise $\rho$} حدود \lr{۰٫۲۰} تا \lr{۰٫۲۴}، \lr{LOO $\rho$} حدود \lr{۰٫۳۱} تا \lr{۰٫۴۰} و \lr{Kendall's W} حدود \lr{۰٫۲۶} تا \lr{۰٫۳۱} است. یعنی نقشه‌ی دو نفر تصادفی فقط کمی شبیه هم است. با این حال \lr{split-half reliability} بین \lr{۰٫۶۲} و \lr{۰٫۷۵} است. پس نقشه‌ی میانگین گروه تکرارپذیر است، حتی اگر نقشه‌ی هر فرد پرنویز باشد؛ مثل یک نظرسنجی که پاسخ تک‌تک افراد پراکنده است ولی میانگین آن پایدار.

سه روش تقریباً به یک اندازه یکسانی دارند. نقشه‌ی \lr{attention flow} با اینکه تقریباً صاف است (بخش~\ref{sec:rollout-flow})، باز هم ترتیب یکسانی بین افراد دارد؛ چون همان تفاوت‌های کوچکش از توجه لایه‌ی اول می‌آید.

\begin{table}[htbp]
\centering
\footnotesize
\caption{همان معیارها برای نقشه‌ی \lr{gradient-weighted rollout} هر عضله در خروجی \lr{E}.}
\label{tab:cons-elec-musc}
\begin{tabular}{lcccccc}
\toprule
نقشه & \lr{pairwise $\rho$} & \lr{LOO $\rho$} & \lr{split-half} & \lr{Kendall's W} & \lr{null 95\%} & \lr{p} \\
\midrule
\lr{E, first-touch, AD} & \lr{۰٫۲۱} & \lr{۰٫۳۴} & \lr{۰٫۶۷} & \lr{۰٫۲۷} & \lr{۰٫۰۴} & \lr{۰٫۰۰۲} \\
\lr{E, first-touch, BR} & \lr{۰٫۲۱} & \lr{۰٫۳۵} & \lr{۰٫۶۸} & \lr{۰٫۲۷} & \lr{۰٫۰۴} & \lr{۰٫۰۰۲} \\
\lr{E, first-touch, FD} & \lr{۰٫۲۱} & \lr{۰٫۳۳} & \lr{۰٫۶۶} & \lr{۰٫۲۷} & \lr{۰٫۰۴} & \lr{۰٫۰۰۲} \\
\lr{E, first-touch, ED} & \lr{۰٫۲۲} & \lr{۰٫۳۷} & \lr{۰٫۶۹} & \lr{۰٫۲۸} & \lr{۰٫۰۴} & \lr{۰٫۰۰۲} \\
\lr{E, first-touch, FDI} & \lr{۰٫۲۱} & \lr{۰٫۳۴} & \lr{۰٫۶۸} & \lr{۰٫۲۷} & \lr{۰٫۰۴} & \lr{۰٫۰۰۲} \\
\lr{E, lift-off, AD} & \lr{۰٫۲۲} & \lr{۰٫۳۵} & \lr{۰٫۶۶} & \lr{۰٫۲۹} & \lr{۰٫۰۴} & \lr{۰٫۰۰۲} \\
\lr{E, lift-off, BR} & \lr{۰٫۲۰} & \lr{۰٫۳۴} & \lr{۰٫۶۵} & \lr{۰٫۲۷} & \lr{۰٫۰۳} & \lr{۰٫۰۰۲} \\
\lr{E, lift-off, FD} & \lr{۰٫۱۸} & \lr{۰٫۳۲} & \lr{۰٫۶۲} & \lr{۰٫۲۵} & \lr{۰٫۰۴} & \lr{۰٫۰۰۲} \\
\lr{E, lift-off, ED} & \lr{۰٫۲۲} & \lr{۰٫۳۸} & \lr{۰٫۶۹} & \lr{۰٫۲۸} & \lr{۰٫۰۴} & \lr{۰٫۰۰۲} \\
\lr{E, lift-off, FDI} & \lr{۰٫۲۰} & \lr{۰٫۳۵} & \lr{۰٫۶۹} & \lr{۰٫۲۶} & \lr{۰٫۰۳} & \lr{۰٫۰۰۲} \\
\bottomrule
\end{tabular}
\end{table}


نقشه‌ی جداگانه‌ی هر عضله (جدول~\ref{tab:cons-elec-musc}) هم همین اندازه یکسانی را دارد و فرق معناداری بین عضله‌ها دیده نمی‌شود. این با بخش~\ref{sec:interp} می‌خواند: در هر شرکت‌کننده، نقشه‌های \lr{gradient-weighted rollout} پنج عضله بسیار شبیه هم‌اند و هیچ عضله‌ای الگوی مکانی خاص خودش را ندارد.

\paragraph{کدام الکترودها در همه مشترک‌اند؟}
جدول~\ref{tab:cons-top} و شکل~\ref{fig:cons_scalp.png} نقشه‌ی گروه را نشان می‌دهند. فقط دو الکترود در سطح گروه به‌طور معنادار سهمی بیشتر از سهم یکنواخت دارند (با تصحیح \lr{max-T} روی ۳۲ الکترود): \lr{C3} و \lr{CP1}. هر دو بالای قشر حسی-حرکتی نیمکره‌ی چپ‌اند که دست راست را کنترل می‌کند. \lr{C3} در ۸ تا ۹ نفر از ۱۲ نفر و \lr{CP1} در ۷ تا ۸ نفر جزو سه الکترود اول‌اند. \lr{C4} در نیمکره‌ی راست سومین الکترود گروه است، اما معنادار نیست. این الگو در هر دو رخداد و هر دو روش \lr{rollout} یکی است.

\begin{table}[htbp]
\centering
\footnotesize
\caption{هشت الکترود با بیشترین سهم میانگین در \lr{gradient-weighted rollout} خروجی \lr{E} هنگام \lr{lift-off}. آزمون $t$ تک‌نمونه‌ای سهم هر الکترود در برابر سهم یکنواخت (\lr{۳٫۱۲۵} درصد) با تصحیح \lr{max-T} روی ۳۲ الکترود.}
\label{tab:cons-top}
\begin{tabular}{lcccc}
\toprule
\lr{electrode} & میانگین گروه (درصد) & \lr{t} در برابر یکنواخت & \lr{p (max-T)} & شمار افراد با آن در سه الکترود اول \\
\midrule
\lr{C3} & \lr{۵٫۶۴} & \lr{۴٫۷۴} & \lr{۰٫۰۲۵} & \lr{9 / 12} \\
\lr{CP1} & \lr{۵٫۰۲} & \lr{۵٫۶۸} & \lr{۰٫۰۰۶} & \lr{7 / 12} \\
\lr{C4} & \lr{۴٫۸۵} & \lr{۲٫۹۶} & \lr{۰٫۳۲۵} & \lr{5 / 12} \\
\lr{PO10} & \lr{۳٫۷۰} & \lr{۱٫۶۶} & \lr{۰٫۹۹۲} & \lr{1 / 12} \\
\lr{T7} & \lr{۳٫۵۳} & \lr{۱٫۰۷} & \lr{۱٫۰۰۰} & \lr{3 / 12} \\
\lr{PO9} & \lr{۳٫۴۰} & \lr{۰٫۹۱} & \lr{۱٫۰۰۰} & \lr{1 / 12} \\
\lr{O1} & \lr{۳٫۳۵} & \lr{۰٫۷۶} & \lr{۱٫۰۰۰} & \lr{2 / 12} \\
\lr{CP2} & \lr{۳٫۲۵} & \lr{۰٫۳۸} & \lr{۱٫۰۰۰} & \lr{2 / 12} \\
\bottomrule
\end{tabular}
\end{table}


\paragraph{شکل نقشه‌ی گروه (شکل~\ref{fig:cons_scalp.png}).} سه جفت نقشه‌ی سر دارد: \lr{rollout} هنگام جدا شدن، \lr{gradient-weighted rollout} هنگام جدا شدن، و \lr{gradient-weighted rollout} هنگام اولین تماس. نقشه‌ی اول هر جفت میانگین سهم الکترودها روی ۱۲ نفر است. نقشه‌ی دوم می‌شمارد هر الکترود در چند نفر جزو سه الکترود اول بوده است (از ۰ تا ۱۲).
\begin{itemize}
\item \textbf{چطور بخوانیم.} نقشه‌ی میانگین ممکن است با چند نفر که سهم خیلی زیادی دارند تیره شود. نقشه‌ی شمارش این مشکل را ندارد: کانون تیره در آن یعنی این الکترود در \emph{تعداد زیادی} از افراد مهم بوده است. اگر دو نقشه یک کانون داشته باشند، آن کانون واقعاً مشترک است.
\item \textbf{چه می‌فهمیم.} در هر سه جفت، قوی‌ترین کانون روی مرکزی چپ (\lr{C3} و \lr{CP1}) است و یک کانون ضعیف‌تر روی مرکزی راست (\lr{C4}). در نقشه‌ی شمارش، مرکزی چپ تا ۸ نفر می‌رسد و بقیه‌ی سر تقریباً سفید است. پس با وجود اینکه نقشه‌ی هر فرد پرنویز است و در بعضی افراد لبه‌های سر هم مهم‌اند، تنها جایی که در اکثر افراد تکرار می‌شود قشر حسی-حرکتی مقابل دست است. این نتیجه به روش (\lr{rollout} یا \lr{gradient-weighted}) و به رخداد (تماس یا جدا شدن) بستگی ندارد. سهم لبه‌ی سر در نقشه‌های میانگین در نقشه‌های شمارش تقریباً دیده نمی‌شود؛ یعنی هر بار در افراد مختلف جای متفاوتی دارد.
\end{itemize}

\fig{cons_scalp.png}{1.0}{نقشه‌ی گروه (میانگین ۱۲ شرکت‌کننده) و تعداد شرکت‌کننده‌هایی که هر الکترود را جزو سه الکترود اول خود دارند، برای \lr{E} با \lr{rollout} و \lr{gradient-weighted rollout}.}

\paragraph{شکل میزان یکسانی (شکل~\ref{fig:cons_electrodes.png}).} هر جفت ستون یک نوع نقشه‌ی الکترود است (خروجی، رخداد، روش). ستون آبی میانگین شباهت دو نفر (\lr{pairwise $\rho$}) و ستون سبز شباهت هر نفر به میانگین بقیه (\lr{LOO $\rho$}) است. خط سیاه کوتاه حدی است که اگر هیچ الگوی مشترکی نبود، فقط از شانس به آن می‌رسیدیم (صدک ۹۵ \lr{null}).
\begin{itemize}
\item \textbf{چطور بخوانیم.} هر چه ستون از خط سیاه بالاتر باشد، شباهت واقعی‌تر است. فاصله‌ی ستون سبز از آبی نشان می‌دهد میانگین گروه چقدر تمیزتر از یک نفر است.
\item \textbf{چه می‌فهمیم.} همه‌ی ۱۲ نوع نقشه پنج تا شش برابر بالاتر از خط شانس‌اند؛ پس شباهت تصادفی نیست. اما مقدار آبی فقط حدود \lr{۰٫۲} است؛ یعنی دو نفر تصادفی فقط کمی شبیه‌اند. سبز تقریباً دو برابر آبی است، پس هر نفر به میانگین گروه شبیه‌تر است تا به یک نفر دیگر. این الگوی «یک نقشه‌ی مشترک به‌علاوه‌ی نویز فردی زیاد» است. ستون‌ها تقریباً هم‌قدند؛ هیچ روش یا رخدادی یکسانی خیلی بهتری نمی‌دهد.
\end{itemize}

\fig{cons_electrodes.png}{1.0}{یکسانی نقشه‌ی الکترودها بین شرکت‌کننده‌ها: \lr{pairwise $\rho$} (آبی) و \lr{LOO $\rho$} (سبز) در برابر صدک ۹۵ \lr{null} (خط سیاه). ستون بالاتر از خط یعنی یکسانی بیشتر از شانس.}

\paragraph{ماتریس شباهت دو‌به‌دو (شکل~\ref{fig:cons_similarity.png}).} هر خانه همبستگی اسپیرمن نقشه‌ی الکترود دو نفر است؛ قرمز یعنی شبیه، آبی یعنی برعکس و سفید یعنی بی‌ربط. قطر همیشه ۱ است.
\begin{itemize}
\item \textbf{چطور بخوانیم.} اگر افراد چند گروه با نقشه‌های متفاوت بودند، در ماتریس بلوک‌های قرمز جدا از هم دیده می‌شد. اگر یک یا دو نفر نقشه‌ی بسیار متفاوتی داشتند، سطر و ستون آن‌ها سفید یا آبی بود.
\item \textbf{چه می‌فهمیم.} تقریباً همه‌ی خانه‌ها قرمز کم‌رنگ‌اند (حدود \lr{۰٫۱} تا \lr{۰٫۴}) و هیچ بلوک یا گروهی دیده نمی‌شود. فقط چند جفت کمی آبی‌اند (مثلاً \lr{P3} و \lr{P4} با \lr{P9}، و \lr{P11} با \lr{P12}). پس افراد زیرگروه ندارند و هیچ‌کس نقشه‌ی کاملاً متفاوتی ندارد. یک الگوی مشترک ضعیف در همه هست و میانگین گرفتن روی گروه کار درستی است.
\end{itemize}

\fig{cons_similarity.png}{0.5}{همبستگی اسپیرمن نقشه‌ی الکترود \lr{E} (\lr{gradient-weighted rollout}، \lr{lift-off}) برای هر جفت شرکت‌کننده.}

\subsection{منحنی‌های زمانی}

\begin{table}[htbp]
\centering
\footnotesize
\caption{یکسانی منحنی‌های زمانی بین شرکت‌کننده‌ها (همبستگی پیرسون؛ توکن خودِ رخداد کنار گذاشته شده). \lr{null}: جابه‌جایی دایره‌ای منحنی هر شرکت‌کننده. ستون آخر: زمان قله نسبت به رخداد، میانگین $\pm$ انحراف معیار بین شرکت‌کننده‌ها.}
\label{tab:cons-time}
\begin{tabular}{lccccccc}
\toprule
نقشه & \lr{pairwise $\rho$} & \lr{LOO $\rho$} & \lr{split-half} & \lr{Kendall's W} & \lr{null 95\%} & \lr{p} & \lr{peak (s)} \\
\midrule
\lr{K, first-touch, rollout} & \lr{۰٫۶۹} & \lr{۰٫۸۱} & \lr{۰٫۹۶} & \lr{۰٫۷۵} & \lr{۰٫۱۲} & \lr{۰٫۰۰۲} & \lr{−۰٫۷۷} $\pm$ \lr{۰٫۸۴} \\
\lr{K, first-touch, gradient-weighted rollout} & \lr{۰٫۵۸} & \lr{۰٫۷۴} & \lr{۰٫۹۵} & \lr{۰٫۶۷} & \lr{۰٫۰۸} & \lr{۰٫۰۰۲} & \lr{−۰٫۴۶} $\pm$ \lr{۰٫۵۳} \\
\lr{K, lift-off, rollout} & \lr{۰٫۶۵} & \lr{۰٫۷۹} & \lr{۰٫۹۶} & \lr{۰٫۶۹} & \lr{۰٫۰۹} & \lr{۰٫۰۰۲} & \lr{−۰٫۵۶} $\pm$ \lr{۰٫۵۹} \\
\lr{K, lift-off, gradient-weighted rollout} & \lr{۰٫۵۴} & \lr{۰٫۷۱} & \lr{۰٫۹۴} & \lr{۰٫۶۷} & \lr{۰٫۰۹} & \lr{۰٫۰۰۲} & \lr{−۰٫۴۴} $\pm$ \lr{۱٫۲۵} \\
\lr{K, first-touch, flow - uniform} & \lr{۰٫۷۱} & \lr{۰٫۸۳} & \lr{۰٫۹۷} & \lr{۰٫۷۳} & \lr{۰٫۱۳} & \lr{۰٫۰۰۲} & -- \\
\lr{K, lift-off, flow - uniform} & \lr{۰٫۶۶} & \lr{۰٫۷۹} & \lr{۰٫۹۶} & \lr{۰٫۶۱} & \lr{۰٫۱۰} & \lr{۰٫۰۰۲} & -- \\
\lr{E, first-touch, rollout} & \lr{۰٫۲۶} & \lr{۰٫۴۸} & \lr{۰٫۸۸} & \lr{۰٫۲۸} & \lr{۰٫۰۸} & \lr{۰٫۰۰۲} & \lr{۰٫۲۳} $\pm$ \lr{۲٫۱۹} \\
\lr{E, first-touch, gradient-weighted rollout} & \lr{۰٫۲۵} & \lr{۰٫۴۵} & \lr{۰٫۸۰} & \lr{۰٫۳۳} & \lr{۰٫۰۷} & \lr{۰٫۰۰۲} & \lr{−۰٫۵۱} $\pm$ \lr{۱٫۲۲} \\
\lr{E, lift-off, rollout} & \lr{۰٫۳۱} & \lr{۰٫۵۳} & \lr{۰٫۹۱} & \lr{۰٫۳۱} & \lr{۰٫۰۸} & \lr{۰٫۰۰۲} & \lr{−۰٫۰۳} $\pm$ \lr{۲٫۳۶} \\
\lr{E, lift-off, gradient-weighted rollout} & \lr{۰٫۲۶} & \lr{۰٫۴۵} & \lr{۰٫۷۹} & \lr{۰٫۳۳} & \lr{۰٫۰۶} & \lr{۰٫۰۰۲} & \lr{−۰٫۸۱} $\pm$ \lr{۱٫۵۶} \\
\lr{E, first-touch, flow - uniform} & \lr{۰٫۱۱} & \lr{۰٫۲۸} & \lr{۰٫۶۹} & \lr{۰٫۱۸} & \lr{۰٫۰۷} & \lr{۰٫۰۲۰} & -- \\
\lr{E, lift-off, flow - uniform} & \lr{۰٫۱۰} & \lr{۰٫۲۷} & \lr{۰٫۶۵} & \lr{۰٫۱۸} & \lr{۰٫۰۷} & \lr{۰٫۰۲۰} & -- \\
\bottomrule
\end{tabular}
\end{table}

\paragraph{شکل منحنی‌های زمانی (شکل~\ref{fig:cons_time.png}).} شش پنل دارد: ردیف بالا \lr{K} و ردیف پایین \lr{E}؛ ستون‌ها از چپ \lr{rollout}، \lr{gradient-weighted rollout} و «\lr{flow} منهای یکنواخت». هر خط خاکستری منحنی یک شرکت‌کننده است و خط رنگی میانگین ۱۲ نفر. شباهت دو‌به‌دو و شباهت هر نفر به بقیه بالای هر پنل نوشته شده است.
\begin{itemize}
\item \textbf{چطور بخوانیم.} اگر خط‌های خاکستری یک دسته‌ی باریک دور خط رنگی بسازند، همه‌ی افراد یک الگوی زمانی دارند. اگر پخش باشند و قله‌هایشان در جاهای مختلف باشد، میانگین چیزی را نشان می‌دهد که در هیچ فردی به آن شکل وجود ندارد.
\item \textbf{چه می‌فهمیم.} در ردیف بالا خط‌های خاکستری دسته‌ی باریکی دارند و یک شکل را تکرار می‌کنند: بالا پیش از رخداد، افت تند درست بعد از آن، کف از ۰ تا \lr{۲٫۵} ثانیه و بالا رفتن در انتهای پنجره. این مطمئن‌ترین نتیجه‌ی تفسیری کل گزارش است: \lr{K} در همه‌ی افراد از حرکت پیش از جدا شدن استفاده می‌کند. در ردیف پایین خط‌ها پخش‌اند، شکل مشترک کمرنگ است و بیشترین ویژگی مشترک بالا رفتن در لبه‌ی پنجره است که به رخداد ربطی ندارد. پس الگوی زمانی \lr{E} را نمی‌توان مثل \lr{K} به همه‌ی افراد تعمیم داد.
\end{itemize}

\fig{cons_time.png}{1.0}{منحنی‌های زمانی در اطراف \lr{lift-off} برای هر ۱۲ شرکت‌کننده (خاکستری) و میانگین آن‌ها (رنگی). بالا: \lr{K}؛ پایین: \lr{E}.}

\paragraph{\lr{K}: بسیار یکسان.} منحنی‌های زمانی \lr{K} در افراد مختلف شکل تقریباً یکسانی دارند (جدول~\ref{tab:cons-time}، شکل~\ref{fig:cons_time.png}): \lr{pairwise r} بین \lr{۰٫۵۴} و \lr{۰٫۷۱}، \lr{LOO r} تا \lr{۰٫۸۳} و \lr{split-half} بالای \lr{۰٫۹۳}. در همه‌ی افراد ورودی‌های پیش از رخداد سهم زیادی دارند و دو ثانیه‌ی بعد از آن کمترین سهم را. قله (در بازه‌ی $\pm 3$ ثانیه) به‌طور میانگین \lr{۰٫۴} تا \lr{۰٫۸} ثانیه پیش از رخداد است، با انحراف معیار \lr{۰٫۵} تا \lr{۱٫۳} ثانیه بین افراد.

\paragraph{\lr{E}: یکسانی کمتر.} منحنی‌های \lr{E} شباهت کمتری دارند (\lr{pairwise r} حدود \lr{۰٫۲۵} تا \lr{۰٫۳۱}) و زمان قله‌شان بین افراد بیشتر جابه‌جا می‌شود (انحراف معیار \lr{۱٫۲} تا \lr{۲٫۴} ثانیه).

\paragraph{\lr{attention flow}.} منحنی «\lr{flow} منهای یکنواخت» برای \lr{K} همان‌قدر یکسان است که \lr{rollout}. اما برای \lr{E} کمترین یکسانی را دارد (\lr{pairwise r} حدود \lr{۰٫۱۰})، که با به سقف رسیدن \lr{attention flow} در جریان \lr{EEG} می‌خواند.

\paragraph{اثر لبه‌ی پنجره.} در همه‌ی افراد و هر دو خروجی، سهم \lr{tokens} نزدیک دو لبه‌ی پنجره‌ی ۸ ثانیه‌ای ($\pm 4$ ثانیه) بالا می‌رود. این ربطی به رخداد ندارد. دلیلش \lr{positional encoding} و این است که \lr{tokens} لبه همسایه‌ی کمتری دارند. چون این الگو در همه مشترک است، بخشی از شباهت منحنی‌ها هم از همین می‌آید. برای همین قله فقط در بازه‌ی $\pm 3$ ثانیه جست‌وجو شده است.

\subsection{\lr{muscle graph}، باندهای \lr{EEG} و مقایسه با استراحت}

\begin{table}[htbp]
\centering
\footnotesize
\caption{یکسانی \lr{muscle graph} (۲۰ درایه‌ی غیرقطری ماتریس ۵×۵) بین شرکت‌کننده‌ها.}
\label{tab:cons-muscle}
\begin{tabular}{lcccccc}
\toprule
نقشه & \lr{pairwise $\rho$} & \lr{LOO $\rho$} & \lr{split-half} & \lr{Kendall's W} & \lr{null 95\%} & \lr{p} \\
\midrule
\lr{K, first-touch, rollout} & \lr{۰٫۳۸} & \lr{۰٫۵۸} & \lr{۰٫۹۰} & \lr{۰٫۴۳} & \lr{۰٫۰۵} & \lr{۰٫۰۰۲} \\
\lr{K, first-touch, gradient-weighted rollout} & \lr{۰٫۵۰} & \lr{۰٫۶۷} & \lr{۰٫۹۳} & \lr{۰٫۵۴} & \lr{۰٫۰۴} & \lr{۰٫۰۰۲} \\
\lr{K, lift-off, rollout} & \lr{۰٫۴۱} & \lr{۰٫۶۲} & \lr{۰٫۹۲} & \lr{۰٫۴۶} & \lr{۰٫۰۵} & \lr{۰٫۰۰۲} \\
\lr{K, lift-off, gradient-weighted rollout} & \lr{۰٫۳۹} & \lr{۰٫۵۷} & \lr{۰٫۸۶} & \lr{۰٫۴۴} & \lr{۰٫۰۵} & \lr{۰٫۰۰۲} \\
\lr{E, first-touch, rollout} & \lr{۰٫۴۷} & \lr{۰٫۶۳} & \lr{۰٫۸۹} & \lr{۰٫۵۱} & \lr{۰٫۰۵} & \lr{۰٫۰۰۲} \\
\lr{E, first-touch, gradient-weighted rollout} & \lr{۰٫۱۸} & \lr{۰٫۲۹} & \lr{۰٫۵۸} & \lr{۰٫۲۵} & \lr{۰٫۰۵} & \lr{۰٫۰۰۲} \\
\lr{E, lift-off, rollout} & \lr{۰٫۴۶} & \lr{۰٫۶۰} & \lr{۰٫۸۶} & \lr{۰٫۵۱} & \lr{۰٫۰۵} & \lr{۰٫۰۰۲} \\
\lr{E, lift-off, gradient-weighted rollout} & \lr{۰٫۱۹} & \lr{۰٫۳۴} & \lr{۰٫۶۳} & \lr{۰٫۲۶} & \lr{۰٫۰۵} & \lr{۰٫۰۰۲} \\
\bottomrule
\end{tabular}
\end{table}

\begin{table}[htbp]
\centering
\footnotesize
\caption{سهم هر باند \lr{EEG} در خروجی \lr{E} هنگام \lr{lift-off} (\lr{gradient × input}). \lr{Kendall's W} روی ترتیب باندها: \lr{۰٫۵۹} (\lr{p} = \lr{<۰٫۰۰۱}).}
\label{tab:cons-bands}
\begin{tabular}{lccc}
\toprule
\lr{band} & سهم میانگین (درصد) & انحراف معیار & شمار افراد با بیشترین سهم \\
\midrule
\lr{13-30 Hz} & \lr{۳۳٫۵} & \lr{۷٫۴} & \lr{9 / 12} \\
\lr{30-40 Hz} & \lr{۲۵٫۲} & \lr{۲٫۷} & \lr{0 / 12} \\
\lr{4-8 Hz} & \lr{۱۶٫۲} & \lr{۴٫۴} & \lr{0 / 12} \\
\lr{8-13 Hz} & \lr{۲۵٫۱} & \lr{۷٫۶} & \lr{3 / 12} \\
\bottomrule
\end{tabular}
\end{table}

\paragraph{شکل باندها (شکل~\ref{fig:cons_bands.png}).} هر ستون یک شرکت‌کننده است و چهار رنگ آن سهم چهار باند \lr{EEG} در \lr{E} هنگام جدا شدن (\lr{gradient × input}). جمع هر ستون ۱ است.
\begin{itemize}
\item \textbf{چطور بخوانیم.} اگر همه‌ی باندها یک‌چهارم سهم داشتند، هر رنگ \lr{۰٫۲۵} ستون را می‌گرفت. رنگی که در بیشتر ستون‌ها بزرگ‌تر از یک‌چهارم است، باندی است که مدل در بیشتر افراد ترجیح می‌دهد.
\item \textbf{چه می‌فهمیم.} بتا (رنگ تیره، پایین هر ستون) در بیشتر افراد بزرگ‌ترین بخش است؛ در \lr{P6} و \lr{P7} به \lr{۰٫۴} تا \lr{۰٫۵} می‌رسد. استثنای روشن \lr{P3} است که آلفا (زرد، بالای ستون) در آن حدود \lr{۰٫۴} سهم دارد. ترتیب باندها از فردی به فرد دیگر کمتر عوض می‌شود تا نقشه‌ی الکترودها؛ یعنی «کدام باند» پایدارتر از «کدام الکترود» است.
\end{itemize}

\fig{cons_bands.png}{0.75}{سهم هر باند \lr{EEG} در خروجی \lr{E} هنگام \lr{lift-off} برای هر شرکت‌کننده.}
\begin{table}[htbp]
\centering
\footnotesize
\caption{یکسانی نقشه‌های $t$ مقایسه‌ی رخداد با استراحت بین شرکت‌کننده‌ها.}
\label{tab:cons-rest}
\begin{tabular}{lcccccc}
\toprule
نقشه & \lr{pairwise $\rho$} & \lr{LOO $\rho$} & \lr{Kendall's W} & \lr{null 95\%} & \lr{p} & الکترودهای معنادار (میانگین) \\
\midrule
\lr{E, first-touch, event - rest t} & \lr{۰٫۰۶} & \lr{۰٫۱۸} & \lr{۰٫۱۴} & \lr{۰٫۰۴} & \lr{۰٫۰۰۶} & \lr{۹٫۹} \\
\lr{E, lift-off, event - rest t} & \lr{۰٫۰۶} & \lr{۰٫۱۷} & \lr{۰٫۱۳} & \lr{۰٫۰۴} & \lr{۰٫۰۱۴} & \lr{۱۱٫۲} \\
\bottomrule
\end{tabular}
\end{table}


\paragraph{\lr{muscle graph}.} درایه‌های غیرقطری \lr{rollout} (یعنی هر عضله چقدر از عضله‌ی دیگر اطلاعات می‌گیرد) در \lr{K} و \lr{E} یکسانی متوسط تا خوبی دارند (جدول~\ref{tab:cons-muscle}، \lr{pairwise $\rho$} حدود \lr{۰٫۳۸} تا \lr{۰٫۵۰}). نسخه‌ی \lr{gradient-weighted} برای \lr{E} یکسانی کمتری دارد (حدود \lr{۰٫۱۸}).

\paragraph{باندهای \lr{EEG}.} باند بتا (۱۳ تا ۳۰ هرتز) با میانگین حدود ۳۴ درصد بیشترین سهم را دارد و در ۹ نفر از ۱۲ نفر اول است؛ در ۳ نفر دیگر آلفا اول است (جدول~\ref{tab:cons-bands}). توافق افراد روی ترتیب چهار باند خوب است (\lr{Kendall's W} حدود \lr{۰٫۵۵} تا \lr{۰٫۶۰}).

\paragraph{مقایسه‌ی رخداد با استراحت.} نقشه‌های $t$ بین افراد تقریباً شبیه هم نیستند (جدول~\ref{tab:cons-rest}، \lr{pairwise $\rho$} حدود \lr{۰٫۰۶})، هرچند هنوز از \lr{null} بیشترند. یعنی اینکه «کدام الکترود در رخداد بیشتر از استراحت سهم دارد» در هر نفر فرق می‌کند. این نقشه‌ها برای ادعای گروهی مناسب نیستند.

\subsection{جمع‌بندی}

\begin{itemize}
\item \textbf{قوی و تکرارپذیر.} الگوی زمانی \lr{K}: همه‌ی افراد به حرکت پیش از رخداد نگاه می‌کنند. همچنین نقشه‌ی گروهی الکترودهای \lr{E} با دو کانون \lr{C3} و \lr{CP1} (\lr{split-half} حدود \lr{۰٫۷}).
\item \textbf{معنادار ولی ضعیف در سطح فرد.} نقشه‌ی الکترود هر شرکت‌کننده (\lr{pairwise $\rho$} حدود \lr{۰٫۲}). ادعا درباره‌ی الکترودها باید درباره‌ی گروه باشد، نه یک نفر.
\item \textbf{قابل اتکا نیست.} نقشه‌های مقایسه‌ی رخداد با استراحت، و «\lr{flow} منهای یکنواخت» برای \lr{E}.
\item \textbf{دیده نشد.} هیچ روشی الگوی خاص هر عضله را به‌طور تکرارپذیر بین افراد نشان نمی‌دهد.
\end{itemize}
